%!TEX root = ../thesis.tex

%jama %hakon -- joint appendix: compile and run instructions

\chapter{Instructions to Compile and Run System}
\label{apx:instructions}

This appendix describes how to reproduce the simulation results presented in this thesis.
The system consists of two parts: the macro-level Simulink/MPC model (developed by H\aa{}kon Rullestad H\aa{}varstein) and the micro-level agent-based MATLAB model (developed by Mohamed Ahmed Jama).

\section{Required Software}
\label{apx:software}

%jama %hakon
\begin{itemize}
  \item MATLAB R2025a or later (earlier R2024 releases may work but are not tested).
  \item Simulink (included with MATLAB).
  \item Model Predictive Control Toolbox for MATLAB \citep{mathworks2025}.
  \item Global Optimization Toolbox (for the GA calibration step).
\end{itemize}

\section{Running the Integrated Simulation}
\label{apx:run}

%jama
To reproduce the results presented in this thesis, follow these steps:

\begin{enumerate}
  \item Open MATLAB (R2025a or later is recommended).
  \item Ensure that the following files are available in the MATLAB working directory:
    \begin{itemize}
      \item \texttt{main.m}
      \item \texttt{run\_micro\_sim.m}
      \item \texttt{update\_agents.m}
      \item \texttt{init\_agents.m}
      \item \texttt{aggregate\_stats.m}
      \item \texttt{sirvhdb\_ode.m}
      \item \texttt{build\_plant\_ga.m}
      \item \texttt{SIRVDB\_model\_parameters.m}
    \end{itemize}
  \item Open the Simulink model \texttt{SIRVHDB\_MPC\_model.slx}.
  \item Run the script \texttt{main.m}.
\end{enumerate}

Executing \texttt{main.m} performs the complete integrated workflow:
\begin{itemize}
  \item the MPC-controlled macro-level simulation is executed,
  \item the transmission-rate trajectories $\beta_1(t)$ and $\beta_2(t)$ are extracted from the Simulink output,
  \item the agent-based simulation is run using these trajectories as inputs,
  \item and the figures reported in the thesis are generated.
\end{itemize}

\section{Simulation Parameters (Micro Model)}
\label{apx:micro_params}

%jama
The main simulation settings used throughout the thesis are summarised below.

\paragraph{Population.}
$N_1 = 500$,\quad $N_2 = 500$.

\paragraph{Simulation horizon.}
$T = 150$~days.

\paragraph{Epidemic parameters.}
$\beta_{\mathrm{nom}} = 0.45$, $\gamma = 1/12$, $\gamma_v = 1/9$, $\gamma_h = 1/25$,
$h = 0.05$, $h_v = 0.02$, $\mu = 0.001$, $\mu_v = 0.0005$, $\mu_h = 0.02$,
$\eta = 0.6$, $\omega_R = 1/110$, $\omega_V = 1/140$,
$\alpha = 0.35$, $r_{\mathrm{inf}} = 2.8$, $p_{\mathrm{commute}} = 0.015$.

\paragraph{MPC parameters (applied to the Simulink model).}
Prediction horizon $N_p = 40$, control horizon $M = 10$, reference $H_{\mathrm{ref}} = 0.8$.

\paragraph{GA calibration.}
Optimised parameters: $\beta_1$, $\gamma$, $h$.
Cost function: MAPE + peak penalty.
Bounds: $\beta_1 \in [0.05, 0.5]$, $\gamma \in [0.05, 0.2]$, $h \in [0.001, 0.1]$.

\section{Simulink Model Parameters (Macro Model)}
\label{apx:simulink_params}

%hakon
The Simulink SIRVHDB model is configured via \texttt{SIRVDB\_model\_parameters.m}.
Key parameters are listed below.

\paragraph{City populations.}
Both cities use normalised compartment values (fractions); the absolute population is not required for ODE integration.

\paragraph{Initial conditions.}
$I_c(0) = 0.01$, $S_c(0) = 0.99$; all other compartments initialised to zero.

\paragraph{Vaccination rates.}
$\nu_1 = 0.001$~day$^{-1}$, $\nu_2 = 0.002$~day$^{-1}$.

\paragraph{Travel matrix.}
\[
T = \begin{pmatrix} -0.2 & 0.2 \\ 0.1 & -0.1 \end{pmatrix}.
\]

\paragraph{Commuting matrix.}
\[
M = \begin{pmatrix} 0.85 & 0.15 \\ 0.10 & 0.90 \end{pmatrix}.
\]

\paragraph{MPC controller.}
Prediction horizon: 40; control horizon: 10; sampling time: 1~day;
output tracking weight: 1.0; control rate weight: 0.2; absolute control weight: 0.05.
Rate limiter on control signal; transport delay: 5~days.

\paragraph{Seasonal forcing.}
Sinusoidal modulation amplitude: 0.4; period: 365~days.

\paragraph{Hospital reference.}
$H_{\mathrm{ref}} = 0.8$ (set in parameter file).

\section{Genetic Algorithm Configuration}
\label{apx:ga_config}

%jama
The genetic algorithm is implemented using MATLAB's \texttt{ga()} function from the Global Optimization Toolbox.
Each candidate solution is a parameter vector
\[
  \theta = (\beta_1,\; \gamma,\; h).
\]
The optimisation variables are constrained by the following bounds:
\[
  \beta_1 \in [0.05,\; 0.5], \qquad
  \gamma \in [0.05,\; 0.2], \qquad
  h \in [0.001,\; 0.1].
\]
The fitness function minimises the mismatch between the ODE and ABM hospital load trajectories using the sMAPE objective with an additional penalty on peak misalignment.
Hospital load is used as the calibration target because it is the primary controlled variable in the MPC formulation and directly represents the healthcare-capacity constraint.
